\section{Operation Modes}\label{operation-modes}

The mode buttons change what the display, knobs and buttons do. The
modes are Voice, Performance (PERF), Step, Load/Save and, with SHIFT,
the Effect page, the pattern generator and the settings menu.

\subsection{Voice Mode}\label{voice-mode}

Voice mode edits sounds and basic patterns.

\begin{itemize}
\tightlist
\item
  \textbf{Display and knobs:} the parameters of the selected synthesis
  page for the active voice.
\item
  \textbf{SEQ buttons:} set and clear the steps of the visible bar.
\item
  \textbf{SELECT buttons:} choose the synthesis page (oscillator,
  amplitude envelope, modulation, FM, transient, filter, LFO, mix).
\item
  \textbf{Voice buttons:} select the active voice and track; with SHIFT,
  mute.
\end{itemize}

\textbf{Example --- change the filter cutoff of voice 3:} press voice
button 3, press the SELECT button for the filter page, turn the first
knob to change \texttt{frq}.

\subsubsection{Step edit with SHIFT}\label{step-edit-with-shift}

Hold SHIFT in Voice mode to show the step parameter page of the
last-touched step. Turn the knobs to edit its specials. A flashing LED
marks the step.

\subsubsection{Morph view}\label{morph-view}

Hold SHIFT (or latch the view with SHIFT + Voice mode) to see and edit
the \textbf{Morph endpoint} of each parameter in Voice mode. A parameter
that cannot be morphed shows its normal value.

\figplaceholder{mode_voice.png}{Voice mode controls.}{Front panel with Voice mode regions highlighted: display, knobs, SEQ buttons, SELECT buttons and voice buttons, each with a one-line label.}

\subsection{Step Mode}\label{step-mode}

Step mode is the inverse of Voice mode with respect to the sequencer:
the SEQ buttons \textbf{select} steps (hold SHIFT to set or clear them)
and the display always shows the step parameter page, so specials and
automation are edited without holding SHIFT. The SELECT buttons show and
edit the step's sub-divisions. Holding several steps edits velocity,
note and condition on all of them.

Hold SHIFT in Step mode to see the Morph endpoints of the track's
length, scale and shuffle. \textbf{SHIFT + Step mode button} enters the
pattern generator.

\verify{The previous LXR described Step mode as selecting steps with the SEQ buttons; confirm the behaviour of SHIFT in Step mode in this firmware.}

\subsection{Performance Mode}\label{performance-mode}

Performance mode is for playing live. Voice editing is not available;
the display shows the PERF pages.

\figplaceholder{mode_perf.png}{Performance mode controls.}{Front panel with the PERF regions highlighted: SEQ buttons as Scene buttons, SELECT buttons as looper divisions, BAR buttons, voice buttons as mutes, display showing mrp fxm 1VM 2VM.}

\subsubsection{Performance menu}\label{performance-menu}

\begin{paramtable}
mrp & Global Morph  & Sets every voice's Morph amount and the Effect Morph at once. 0 to 255. \\
fxm & Effect Morph  & Morph amount of the Scene's Effect. \\
1VM -- 6VM & Voice Morph & Morph amount of one voice. \\
\end{paramtable}

The global Morph control is never recorded. The individual voice and
Effect Morph cells are.

\subsubsection{Scene change}\label{scene-change}

\begin{itemize}
\tightlist
\item
  \textbf{SEQ button:} switch to that Scene. The change happens at the
  \textbf{Scene change division} \texttt{scd} (immediately if the
  sequencer is stopped).
\item
  \textbf{SHIFT + SEQ button:} \emph{queue} a change at the
  \textbf{queue division} \texttt{qcd}; a new queue replaces a pending
  one. The queued Scene's LED flashes at double rate.
\item
  \textbf{Reset behaviour:} \texttt{pcr} (pattern change reset, on the
  settings menu) decides whether a direct change restarts every track
  from step 1. If it is off, the new Scene is \emph{realigned} to the
  running master clock and the effect sequence stays aligned too. Queued
  changes reset only when the queue action \texttt{qac} includes a
  reset.
\end{itemize}

\begin{paramtable}
scd & Scene change division & 0 to 9: instant, 1/8 note, 1/4, 1/2, 1 bar, 2 bars, 3 bars, 4 bars, end of shortest track, end of longest track. \\
qcd & Queue division        & 0 to 17: instant and a longer range of divisions from 1/8 note up to 16 bars, and end of shortest or longest track. \\
qac & Queue action         & What happens when a queued change fires. Includes reset (\texttt{res}), unmute-and-reset (\texttt{u+r}), morph-reset variants (\texttt{mrp}, \texttt{o+r}, \texttt{mmr}), reload (\texttt{rld}) and \texttt{r+r}. 0 to 8. \\
pcr & Pattern change reset & On: direct and chained Scene changes restart the pattern. \\
\end{paramtable}

\verify{Division names shown on screen for scd and qcd, and the full list of nine qac values. The ``1/8 note, 1/4, 1/2, bar'' labels here are derived from the stored step counts.}

\subsubsection{Scene chains}\label{scene-chains}

Each Scene can chain automatically to another. Open the chain editor
with \textbf{SHIFT + Play + SEQ} (hold SHIFT, tap Play, then press the
Scene's SEQ button). The editor, a latched menu, shows
\texttt{cn1\ nx1\ cn2\ nx2}: two \emph{conditions}, each with a
\emph{next} Scene.

\begin{itemize}
\tightlist
\item
  \textbf{Conditions} can be: always, after N repeats of the pattern,
  after N bars, a probability at a repeat or bar end, a bar number
  within a group of four or eight, or after a particular previous Scene.
  Conditions can also test whether fill or loop is active.
\item
  \textbf{Next} can be: the same Scene (a repeat), the next or previous
  Scene, a specific Scene, a random Scene in a range, a random next or
  previous few, or \textbf{stop}.
\item
  Condition 1 has priority; two probabilities that add to 100 or more
  are drawn proportionally.
\item
  Chain conditions are not evaluated while a Bank is loading.
\end{itemize}

\figplaceholder{scene_chain_example.png}{Scene chain example.}{Timeline diagram: Scene 1 plays 3 times, then Scene 2 plays once with probability 50 percent, otherwise Scene 1 repeats. Show the chain editor display below the timeline.}

\subsubsection{Looper}\label{looper}

Hold a SELECT button to loop the playing tracks. The eight buttons set
the division: 1 bar, 1/2, 1/4, 1/8, 1/16, 1/32, 1/64, 1/128. The last
pressed wins.

\begin{itemize}
\tightlist
\item
  Hold a \textbf{voice button} first, then press SELECT, to loop only
  those tracks.
\item
  With the loop running, a voice button adds that track to the loop
  (with a fresh loop start); releasing it drops the track, which resumes
  where it would have been.
\item
  \textbf{BAR 1 / BAR 2} while looping apply a momentary
  \textbf{triplet} (2/3) or \textbf{dotted} (3/2) length.
\item
  Releasing every loop button returns the sequencer to where it would
  have been without looping.
\item
  A track that is not \emph{loopable} (see the STEP page) is ignored by
  the looper.
\item
  While the loop is engaged the SELECT, BAR and voice buttons keep their
  looper meaning; live recording keeps running.
\end{itemize}

\subsubsection{Trigger overlay (BAR
buttons)}\label{trigger-overlay-bar-buttons}

When not looping, \textbf{BAR 1} and \textbf{BAR 2} toggle the
\texttt{\textless{}} and \texttt{\textgreater{}} trigger overlays (one
at a time). With an overlay on, the voice buttons \textbf{fire the
overlay trigger} instead of selecting or muting; release is the
note-off. SHIFT + voice button still mutes.

Each overlay is configured on the record settings page (SHIFT + Rec):

\begin{paramtable}
pb< / pb> & Overlay trigger type & Single hits at velocity 127, 100, 64 or 32 (\texttt{vmx v10 v64 v32}); two flams (\texttt{flm fls}); or a held roll at \texttt{1/8 8tr 16 16t 32 32t 64 64t 128}. \\
pn< / pn> & Overlay note & Pitch of the trigger. \texttt{def} is the default pitch. \\
pv< / pv> & Roll rhythm & Velocity rhythm of the roll repeats (the same list as the step \texttt{rtm}). \\
\end{paramtable}

Rolls repeat until you release the button. Tracks have independent
rolls, so several can run at once. The overlay also plays from MIDI
notes (see the MIDI chapter), and rolls keep running while the sequencer
is stopped.

\subsubsection{Manual re-align and per-track
Scenes}\label{manual-re-align-and-per-track-scenes}

See \emph{Per-track Scene playback} in the Sequencer chapter.

\subsection{Live Recording}\label{live-recording}

Live record writes what you play and turn onto \textbf{one recording
track} of one Scene.

{\def\LTcaptype{none} % do not increment counter
\begin{longtable}[]{@{}
  >{\raggedright\arraybackslash}p{(\linewidth - 4\tabcolsep) * \real{0.3333}}
  >{\raggedright\arraybackslash}p{(\linewidth - 4\tabcolsep) * \real{0.3333}}
  >{\raggedright\arraybackslash}p{(\linewidth - 4\tabcolsep) * \real{0.3333}}@{}}
\toprule\noalign{}
\begin{minipage}[b]{\linewidth}\raggedright
Gesture
\end{minipage} & \begin{minipage}[b]{\linewidth}\raggedright
Mode
\end{minipage} & \begin{minipage}[b]{\linewidth}\raggedright
Behaviour
\end{minipage} \\
\midrule\noalign{}
\endhead
\bottomrule\noalign{}
\endlastfoot
Tap \textbf{Rec} & LATCH & Recording toggles on release. The first
adjustment of a parameter \emph{latches} it: during the take the
parameter owns its automation on the recording track. \\
Hold \textbf{Rec} (more than 300 ms) & OVERDUB & Only parameters you
adjust are written, only while the button is held. Nothing latches or
clears. \\
\textbf{SHIFT + Clear} while recording & ERASE & Live erase (see
below). \\
\end{longtable}
}

Arming while stopped is allowed; recording begins when playback does.
Any mode change of the recorder, a change of recording track or any
Scene change starts a new take.

\subsubsection{What is recorded}\label{what-is-recorded}

\begin{itemize}
\tightlist
\item
  Triggers from the BAR overlay and from MIDI notes (single hits, flams
  and held rolls). A held roll is stored as a roll step at its true rate
  and length; there is no slow/fast split.
\item
  Parameter moves on the Voice pages, per-voice Scene mix cells, Effect
  parameters and the Effect Morph cell, track length, scale and shuffle
  (when the viewed Scene is the recording Scene), and the PERF voice and
  Effect Morph cells.
\item
  \textbf{Not recorded:} the global Morph \texttt{mrp}, settings and
  Load/Save pages, step editors, effect sequence locks, LFO and velocity
  targets, the effect \texttt{typ}, and Morph-view edits.
\end{itemize}

\subsubsection{Timing}\label{timing}

Hits are placed where they were played, \emph{including} the time the
note spent waiting to be handled. A hit that lands slightly ahead of the
playhead is stored on the next step with an echo mark so the step does
not sound twice. The sequence follows the \textbf{quantise} setting.

\subsubsection{Record settings (SHIFT +
Rec)}\label{record-settings-shift-rec}

\begin{paramtable}
qnt & Quantise & \texttt{off}, \texttt{1/8}, \texttt{1/16}, \texttt{1/32}, \texttt{1/64}, or \texttt{trk} (snap to the track's own step grid). Applied while recording; it cannot be applied afterwards. \\
rtk & Recording track & The track that receives the recording, 1 to 7. Follows the active track; SHIFT + Rec can override it. \\
\end{paramtable}

\verify{Whether the quantise values appear as shown on the display and where qnt and rtk sit: they share a page with pv< pv> on the second SHIFT+REC sub-page.}

\subsubsection{Live erase}\label{live-erase}

Hold \textbf{SHIFT + Clear} while recording to erase:

\begin{itemize}
\tightlist
\item
  \textbf{Turn a knob} over an automatable parameter: that parameter's
  automation is erased from the recording track for every step the
  sequence passes while erase is on.
\item
  \textbf{Hold a SEQ button} in Voice or Step mode: that step's triggers
  and specials on the active track are removed and not played
  (automation still plays).
\item
  \textbf{Hold a voice button:} that track's triggers, specials,
  conditions and automation are removed and not played.
\item
  On the Effect page, holding a SEQ button removes the locks of each
  Effect step the effect playhead reaches.
\end{itemize}

\subsection{Morph}\label{morph}

The Morph function transitions from the Normal value of each parameter
to its Morph value.

\begin{itemize}
\tightlist
\item
  \textbf{Where Morph values live:} every morphable parameter has a
  Morph value. Hold SHIFT in the Voice, Step or Effect pages to edit the
  Morph endpoint.
\item
  \textbf{Where Morph amounts come from:} the PERF menu (\texttt{mrp},
  \texttt{fxm}, \texttt{1VM} to \texttt{6VM}), MIDI CC1, step
  automation, velocity, and LFO. CC1 on the global channel sets the
  global Morph; CC1 on a voice's channel sets that voice's Morph and is
  recorded.
\item
  Morph \textbf{never blends} selector parameters such as oscillator
  type, run mode, chain type or the Alkaline waveform; those step
  instead.
\item
  Global Morph also sets the Effect Morph. A global Morph move wins over
  the Effect's Morph lane lock and over \texttt{fxm} automation until
  the next Morph automation.
\end{itemize}

\subsection{Effect Page}\label{effect-page}

Press \textbf{SHIFT + Performance mode}. The Effect page edits the
Scene's Effect and its 16-step effect sequence. See the Mixer and
Effects chapter.

\subsection{Pattern Generator}\label{pattern-generator}

\draftnote{The firmware has an Euclidean generator and a SOM generator, but the engineering notes
describe the generator page as a proposal (``one Generator page, pick a type, Off/Overdub/Replace'')
and report that the SOM generator is not currently wired in. Do not document generators in the
release manual until the generator page exists. The earlier LXR's pattern generator
(\texttt{len stp rot ssr}) is the model for the Euclidean part.}

\subsection{Load/Save Mode}\label{loadsave-mode}

See the Storage chapter.

\subsection{Settings Menu}\label{settings-menu}

Press \textbf{SHIFT + Load/Save}. The settings menu holds the global
settings. Its pages are listed below. Settings are stored in
\texttt{settings.cfg} on the card, which is read at boot.

\begin{paramtable}
snc & Sync source & \texttt{off usb din pls aut}: external clock source. \texttt{aut} picks the highest priority source present (jack pulse over DIN over USB over internal). \\
ch  & Global MIDI channel & \texttt{off} or 1 to 16. \\
txf / rxf & MIDI send / receive filter & Which message classes are sent / acted on: notes (N), realtime (R), control (C), program change (P). \\
mrt & MIDI routing & Pass-through routing between DIN and USB. \\
pcr & Pattern change reset & See the Performance mode section. \\
scd qcd qac & Scene change timing & See the Performance mode section. \\
cmp cam ctm csc & Bus compressor & Last page of the menu; see the Mixer and Effects chapter. \\
\end{paramtable}

\verify{The settings menu has further entries that the project documents do not list (tempo, display, interpolation, fetch behaviour, dither and so on). Take the full list from the firmware's Global page.}

\tip{A card with no MIDI filter keys loads with send, receive and routing all set to \texttt{off}:
the LXR will not consume, send or forward any MIDI until you set \texttt{rxf} and \texttt{txf} to
\texttt{all}.}
